\subsection{An all-prime fixed-subgroup comparison for the tripod}

\begin{statement}{A.34}{Proved}{The tripod over an abelian closure, for every cusp packet}
Let $k$ be a number field in $\overline{\Q}$, let
$V=\mathbf P^1_k\setminus\{0,1,\infty\}$, and fix a prime $\ell$.
Write $F_\ell=\pi_1(V_{\bar k},\overrightarrow{01})^{(\ell)}$
and $I_0=\langle x\rangle\simeq\Z_\ell(1)$ for its actual inertia
at zero. Let $s_b:G_k\longrightarrow F_\ell\rtimes G_k$ be any
continuous section in this cusp packet, with its representative
chosen in $I_0\rtimes G_k$. Put
$\Delta_k=\operatorname{Gal}(\bar k/k^{\rm ab})$.
For every open subgroup $\Delta_0\subset\Delta_k$ we have
\[
                   F_\ell^{s_b(\Delta_0)}=I_0.        \tag{A.34.1}
\]
The same assertion holds at either of the other two cusps after
transporting the base point and inertia by the corresponding
automorphism of $V$. It includes $\ell=2,3$, arbitrary completed
tangent parameters, and packet representatives whose Kummer
character vanishes on $\Delta_0$.

The elementary algebra behind this assertion is the following
unbounded calculation. On the free Lie algebra
$L=\mathbb L_{\Q_\ell}(X,Y)$ set
\[
 F_3=[X,[X,Y]]+[Y,[X,Y]],\qquad
 dX=0,\qquad dY=[Y,F_3].                               \tag{A.34.2}
\]
Then $\ker(d:L\to L)=\Q_\ell X$.
The comparison (A.34.1) is a statement about the actual free
pro-$\ell$ group, not just its rational Lie algebra. It supplies
a target fixed-subgroup calculation for Task A.4; it does not
assert that a cover's maximal pro-$\ell$ quotient embeds in
that of this target.
\end{statement}
\begin{proof}
We first prove the algebraic claim over any field of
characteristic zero. Embed $L$ in the free associative algebra
$T=\Q_\ell\langle X,Y\rangle$. On words of a fixed length use
lexicographic order with $X>Y$. Expanding (A.34.2) gives
\[
\begin{split}
 dY={}&-X^2Y^2+2XYXY-2YXYX+Y^2X^2\\
      &+3Y^2XY-3YXY^2+XY^3-Y^3X .
\end{split}
                                                        \tag{A.34.3}
\]
Its leading word is $-X^2Y^2$. Every word containing $Y$ has a
unique expression $w=X^aYv$ where the displayed $Y$ is its first.
In the Leibniz expansion of $dw$, differentiation at that first
$Y$ has leading word
\[
                         -X^{a+2}Y^2v.               \tag{A.34.4}
\]
Differentiating any later $Y$ leaves a $Y$ in position $a+1$,
whereas (A.34.4) has an $X$ there. All such terms are smaller.
The leading words (A.34.4), as $w$ varies in one degree, are
distinct and preserve the lexicographic order: different values
of $a$ are compared at the first $Y$, and equal values are
compared in $v$. Thus a nonzero linear combination of words
containing $Y$ cannot have derivative zero. Consequently
\[
                    \ker(d:T\to T)=\Q_\ell[X].
\]
Its intersection with $L$ is $\Q_\ell X$. For example this last
fact follows from the primitive coproduct: $X^n$ is not
primitive for $n>1$ in characteristic zero. This proves the
kernel assertion in every degree. The same word argument, or
comparison of the first and last letters, gives
\[
       \ker(\operatorname{ad}X:L\to L)=\Q_\ell X.    \tag{A.34.5}
\]

We next obtain an actual Galois element with first derivation a
nonzero multiple of $d$. The arithmetic input is the nonzero
degree-three Soul\'e class in the actual outer Johnson filtration.
Precisely, with
\[
 I_m^\ell G_\Q=\ker\{G_\Q\longrightarrow
             \operatorname{Out}(F_\ell/\gamma_{m+1}F_\ell)\},
\]
the rational vector space
$(I_3^\ell G_\Q/I_4^\ell G_\Q)\otimes\Q_\ell$ is nonzero.
This is the $m=1$ case of
\cite[Remark 8.6, together with Theorem 8.4 and
the filtration comparison in Proposition 8.2]{A_HainMatsumoto2003}.
That source states the nonvanishing for every prime; it is not
an assertion only for regular primes or only for $\ell>3$.
The comparison with the based tangential action is specified
in its Section 8.5. Only this degree-three nonvanishing is
used here, not the freeness of the full arithmetic Lie algebra.

Choose $h\in I_3^\ell G_\Q$ with nonzero rational degree-three
class. For the unit tangential section its genuine action is
\[
 \rho_0(g)x=x^{\chi_\ell(g)},\qquad
 \rho_0(g)y=f_g^{-1}y^{\chi_\ell(g)}f_g,
                  \qquad f_g\in F_\ell'.             \tag{A.34.6}
\]
Here the interval path from zero to one fixes the normalization;
in particular $f_g$ has zero abelianization and satisfies the
two-cycle identity $f_g(x,y)f_g(y,x)=1$.
We check that passing from the outer filtration to (A.34.6)
introduces no lower inner term. Since $\chi_\ell(h)=1$ and
$f_h\in\gamma_2F_\ell$, this action is already trivial modulo
$\gamma_3F_\ell$. If its possible degree-two derivation were
inner, its inner element's degree-one part would have zero
bracket with both $X,Y$, and hence be zero. Its degree-two
part would commute with $X$, and hence be zero by (A.34.5).
Therefore the degree-two derivation vanishes. The degree-two
part of $f_h$ is a multiple of $[X,Y]$; its bracket with $Y$
is zero only for the zero multiple. Thus
$f_h\in\gamma_3F_\ell$.

Write its degree-three part as
$a[X,[X,Y]]+b[Y,[X,Y]]$. Interchanging $X,Y$ sends these two
basis elements to the negatives of one another. The two-cycle
identity consequently gives $a=b$. Nonzero outer degree three
forces $a\ne0$, since a zero degree-three part would make the
actual action trivial modulo $\gamma_5F_\ell$. The first
derivation of $\rho_0(h)$ is therefore $a d$, up to the
irrelevant simultaneous sign convention for group commutators.

Take a positive integer power of $h$ to put it in the open
subgroup $G_k$; its first derivation is multiplied by that
integer and stays nonzero. Choose $g\in G_k$ with
$\chi_\ell(g)^3\ne1$, which is possible because the cyclotomic
image of a number field is open in $\Z_\ell^\times$. Conjugation
by $\rho_0(g)$ multiplies a degree-three derivation by
$\chi_\ell(g)^3$: on homogeneous degree $n$ the underlying
associated graded action is multiplication by
$\chi_\ell(g)^n$. Hence the actual commutator
\[
                         e=[g,h]\in\Delta_k
\]
has first derivation
$(\chi_\ell(g)^3-1)a d\ne0$. It still belongs to
$I_3^\ell G_\Q$. Taking a further positive integer power puts
$e$ in the specified open subgroup $\Delta_0$, again without
killing this rational leading derivation. This establishes
the required arithmetic element in every $\Delta_0$ by actual
group operations.

Every representative of the packet has the form
$s_b(g)=x^{\kappa(g)}s_0(g)$ for an integral continuous
cocycle $\kappa:G_k\to\Z_\ell(1)$. Thus
\[
                 \rho_b(g)=\operatorname{Inn}(x^{\kappa(g)})
                                  \rho_0(g).          \tag{A.34.7}
\]
For the element $e$ just constructed, there are two cases.
If $\kappa(e)\ne0$, the first derivation of $\rho_b(e)$ is
the degree-one derivation $\kappa(e)\operatorname{ad}X$.
If $\kappa(e)=0$, its first derivation is the nonzero multiple
of $d$ already obtained. Both kernels on $L$ are exactly
$\Q_\ell X$. This treats the zero-Kummer packet as well as all
other packets; no generic tangent condition is imposed.

Finally, the closed lower central filtration of the free
pro-$\ell$ group is separated, with
\[
 \operatorname{gr}_\gamma F_\ell
                  \simeq\mathbb L_{\Z_\ell}(X,Y).
\]
One can use Hall commutators to verify this integral comparison:
in each nilpotent quotient their ordered exponent coordinates
are in $\Z_\ell$, and their successive commutator coordinates
are the free Lie basis. Separation also follows directly by
detecting a nonidentity element in a finite $\ell$-group,
which is nilpotent. In particular every nonzero graded class
remains nonzero after tensoring with $\Q_\ell$.

If $u\in F_\ell$ is fixed by $\rho_b(e)$, its first nonzero
lower central class must be killed by the first derivation.
For a class of degree $n$, this follows by taking the
degree $n+j$ component of $\rho_b(e)u\,u^{-1}$, where $j$ is
the first derivation's degree. The two kernel calculations
show that $n=1$ and the abelianization of $u$ is $cX$.
Its coefficient $c$ lies in $\Z_\ell$ because the integral
abelianization is $\Z_\ell X\oplus\Z_\ell Y$. The element
$ux^{-c}$ is still fixed and belongs to $\gamma_2F_\ell$;
the same argument forces it to be one. We conclude
$F_\ell^{\rho_b(e)}=\langle x\rangle$. Since all of
$\Delta_0$ has cyclotomic character one, it fixes $x$;
intersecting fixed groups proves (A.34.1).

The program \texttt{code/a\_tripod\_fixed.py} independently
expands (A.34.3) over the integers and verifies the unique,
order-preserving leading word on every word through degree
twelve. These finite checks support the displayed identities;
the unbounded proof is the word argument above.
\end{proof}
